\ulohahead{specializace UI-SU}{Statistické testy: t-test a chí-kvadrát}
\begin{zadani}
\begin{enumerate}[label=(\alph*)]
  \item \bod{1} Co je statistický test, nulová hypotéza, hladina významnosti a p-hodnota? Popište jednovýběrový a dvouvýběrový t-test.
  \item \bod{2} Popište chí-kvadrát test dobré shody: testovou statistiku, stupně volnosti a rozhodovací pravidlo.
  \item \bod{3} Co je párový t-test a kdy ho použít? Jak souvisí rozhodnutí testu s chybami 1.\ a 2.\ druhu?
\end{enumerate}
\end{zadani}

\begin{reseni}
\textbf{(a)} \emph{Test} rozhoduje mezi \emph{nulovou hypotézou} $H_0$ a alternativou na základě dat. \emph{Hladina významnosti} $\alpha$ = přípustná pravděpodobnost chyby 1.\ druhu; \emph{p-hodnota} = pravděpodobnost dat aspoň tak extrémních jako pozorovaná, platí-li $H_0$; zamítáme, je-li $p<\alpha$. \emph{Jednovýběrový} t-test: $H_0:\mu=\mu_0$, statistika $t=\frac{\bar x-\mu_0}{s/\sqrt n}$, která má za platnosti $H_0$ Studentovo $t$-rozdělení s $n-1$ stupni volnosti; zamítáme, padne-li $t$ do kritického oboru (resp.\ $p<\alpha$). \emph{Dvouvýběrový}: $H_0:\mu_1=\mu_2$, statistika $t=\frac{\bar x-\bar y}{\sqrt{s_X^2/n+s_Y^2/m}}$ (Welchova varianta při různých rozptylech).

\textbf{(b)} \emph{Chí-kvadrát test dobré shody} porovná pozorované četnosti $O_i$ s očekávanými $E_i$ podle modelu: $\chi^2=\sum_i\frac{(O_i-E_i)^2}{E_i}$. \emph{Stupně volnosti} $=$ (počet kategorií $-1-$ počet odhadovaných parametru). $H_0$ (data odpovídají modelu) zamítneme, je-li $\chi^2$ větší než kritická hodnota $\chi^2_{1-\alpha}$ pro dané df.

\textbf{(c)} \emph{Párový} t-test se použije, jsou-li data \emph{spárovaná} (např.\ měření před/po na týchž jedincích); testuje, zda je průměr rozdílu nulový. \emph{Chyba 1.\ druhu} = zamítnout platné $H_0$ (řízena $\alpha$); \emph{chyba 2.\ druhu} = nezamítnout neplatné $H_0$ (řízena silou testu $1-\beta$, roste s velikostí vzorku a efektu).
\end{reseni}

\begin{jadro}
Test: $H_0$, testová statistika, p-hodnota; zamítni při $p<\alpha$. Jednovýběrový t-test $t=\frac{\bar x-\mu_0}{s/\sqrt n}$. Chí-kvadrát dobré shody $\chi^2=\sum_i\frac{(O_i-E_i)^2}{E_i}$, df $=$ kategorie $-1-$ parametry.
\end{jadro}

\kalibrace{Statistické testy (~22\,\% dohromady) chtějí převážně definice + dosazení do statistiky. Stačí 2-bodová záloha: hypotéza, statistika, rozhodnutí.}
\past{p-hodnota není pravděpodobnost, že $H_0$ platí. Stupně volnosti u chí-kvadrát nejsou počet kategorií, ale po odečtení vazeb/parametru.}
